\documentclass[12pt,a4paper]{article}

% Hỗ trợ tiếng Việt với pdfLaTeX
\usepackage[utf8]{inputenc}
\usepackage[T5]{fontenc}
\usepackage[vietnamese]{babel}
\usepackage{lmodern}

% Toán học và định dạng
\usepackage{amsmath,amssymb}
\usepackage{geometry}
\usepackage{booktabs}
\usepackage{tabularx}
\usepackage{array}
\usepackage{enumitem}
\usepackage{microtype}
\usepackage{xcolor}
\usepackage[hidelinks]{hyperref}

\geometry{
    top=2.2cm,
    bottom=2.2cm,
    left=2.3cm,
    right=2.3cm
}

\setlength{\parindent}{0pt}
\setlength{\parskip}{0.6em}
\renewcommand{\arraystretch}{1.25}

\title{\textbf{Đánh giá nền tảng ANN, CNN, AlexNet\\và lộ trình học Reinforcement Learning}}
\author{}
\date{}

\begin{document}

\maketitle
\tableofcontents
\newpage

\section{Kết luận}

\textbf{Bạn đã đủ nền tảng để bắt đầu học RL ở mức nhập môn}, đặc biệt là \textbf{tabular Q-Learning}. Tuy nhiên, ANN/CNN/AlexNet \textbf{không phải các kiến thức nền trực tiếp quan trọng nhất của RL}.

Bạn chưa nên nhảy thẳng vào \textbf{Deep Reinforcement Learning như DQN, PPO hoặc Actor--Critic}. Trước tiên cần bổ sung một số kiến thức về xác suất, quá trình Markov và các thuật toán RL cơ bản.

Theo lộ trình hiện tại, bạn đã hoàn thành phần lớn \textbf{ANN/Backpropagation}, \textbf{CNN} và \textbf{AlexNet 1}; trong khi các nội dung \textbf{Q-Learning} vẫn chưa học.

\section{Những gì bạn đã học có ích thế nào cho RL?}

\subsection{ANN và Backpropagation: rất hữu ích cho Deep RL}

Bạn đã học:

\begin{itemize}[leftmargin=1.5em]
    \item Cấu trúc mạng neural.
    \item Forward propagation và backpropagation.
    \item Gradient descent.
    \item Softmax.
    \item Vanishing gradient.
    \item Underfitting và overfitting.
\end{itemize}

Đây là nền tảng tốt để sau này hiểu:

\begin{itemize}[leftmargin=1.5em]
    \item Mạng Q trong DQN.
    \item Policy network.
    \item Value network.
    \item Actor và Critic.
    \item Cách tối ưu loss trong Deep RL.
\end{itemize}

\subsection{CNN: chỉ cần trong một số bài toán RL}

CNN hữu ích khi trạng thái của môi trường là hình ảnh, chẳng hạn:

\begin{itemize}[leftmargin=1.5em]
    \item Chơi game Atari từ khung hình.
    \item Điều khiển robot bằng camera.
    \item Xe tự hành.
    \item Điều hướng dựa trên hình ảnh.
\end{itemize}

Với các môi trường ban đầu như FrozenLake, Taxi, CartPole hoặc MountainCar, bạn \textbf{không cần CNN}.

\subsection{AlexNet: không phải điều kiện tiên quyết}

AlexNet giúp bạn hiểu một kiến trúc CNN cụ thể, nhưng không giúp nhiều cho RL nhập môn. Bạn không cần hoàn thành AlexNet 2 hoặc phần Matlab AlexNet trước khi học RL.

\section{Các phần trong lộ trình hiện tại: phần nào cần cho RL?}

\begin{table}[htbp]
\centering
\small
\begin{tabularx}{\textwidth}{
    >{\raggedright\arraybackslash}p{4.1cm}
    >{\centering\arraybackslash}p{3.0cm}
    >{\raggedright\arraybackslash}X
}
\toprule
\textbf{Chủ đề trong sheet} &
\textbf{Có cần trước RL không?} &
\textbf{Nhận xét} \\
\midrule
Supervised và Unsupervised Learning &
Nên biết &
Cần hiểu RL khác supervised learning ở điểm nào. \\

Linear Regression &
Nên học &
Hữu ích để hiểu function approximation, prediction và loss. \\

KNN &
Không bắt buộc &
Có thể bỏ qua trước khi học RL. \\

K-Means &
Không bắt buộc &
Không liên quan trực tiếp đến RL cơ bản. \\

SVM &
Không bắt buộc &
Không phải nền tảng của RL. \\

Perceptron &
Không cần học lại &
ANN bạn đã học đã bao hàm ý tưởng này. \\

AlexNet 2 &
Không bắt buộc &
Chỉ cần khi quan tâm đến CNN chuyên sâu. \\

Q-Learning &
Bắt buộc &
Đây là điểm bắt đầu trực tiếp của RL. \\

RNN/LSTM &
Chưa cần &
Chỉ cần sau này với trạng thái không quan sát đầy đủ hoặc dữ liệu chuỗi. \\

Genetic Algorithm &
Không bắt buộc &
Là hướng tối ưu khác, không phải nền tảng RL. \\
\bottomrule
\end{tabularx}
\end{table}

Như vậy, bạn \textbf{không cần quay lại học toàn bộ KNN, K-Means và SVM rồi mới bắt đầu RL}.

Trong các nội dung chưa học của sheet, phần nên bổ sung nhất là:

\begin{enumerate}[leftmargin=1.8em]
    \item Khái niệm supervised, unsupervised và reinforcement learning.
    \item Linear regression ở mức cơ bản.
    \item Các nội dung Q-Learning.
\end{enumerate}

\section{Những kiến thức còn thiếu nhưng sheet chưa đề cập đầy đủ}

\subsection{Xác suất và kỳ vọng --- bắt buộc}

Bạn cần hiểu:

\begin{itemize}[leftmargin=1.5em]
    \item Biến ngẫu nhiên.
    \item Xác suất có điều kiện.
    \item Phân phối xác suất.
    \item Kỳ vọng $\mathbb{E}[X]$.
    \item Trung bình có trọng số theo xác suất.
    \item Sampling.
    \item Variance ở mức cơ bản.
\end{itemize}

Trong RL, phần thưởng và trạng thái tiếp theo thường mang tính ngẫu nhiên. Giá trị của một trạng thái thực chất là \textbf{kỳ vọng của tổng phần thưởng trong tương lai}:

\begin{equation}
V^\pi(s)
=
\mathbb{E}_{\pi}
\left[
\sum_{t=0}^{\infty}
\gamma^t R_{t+1}
\;\middle|\;
S_0=s
\right].
\end{equation}

Nếu chưa hiểu kỳ vọng, công thức value function và Bellman equation sẽ khá khó tiếp cận.

\subsection{Markov Decision Process --- bắt buộc}

Đây là nền tảng quan trọng nhất của RL.

Bạn cần phân biệt rõ:

\begin{itemize}[leftmargin=1.5em]
    \item \textbf{Agent}: tác nhân học.
    \item \textbf{Environment}: môi trường.
    \item \textbf{State $s$}: trạng thái.
    \item \textbf{Action $a$}: hành động.
    \item \textbf{Reward $r$}: phần thưởng.
    \item \textbf{Transition}: cách môi trường chuyển trạng thái.
    \item \textbf{Policy $\pi(a\mid s)$}: chiến lược chọn hành động.
    \item \textbf{Episode}: một lượt tương tác hoàn chỉnh.
    \item \textbf{Terminal state}: trạng thái kết thúc.
    \item \textbf{Discount factor $\gamma$}: mức coi trọng phần thưởng tương lai.
    \item \textbf{Return $G_t$}: tổng phần thưởng tích lũy.
\end{itemize}

Bạn cũng cần hiểu \textbf{Markov property}:

\begin{quote}
Khi biết trạng thái hiện tại, tương lai không phụ thuộc trực tiếp vào toàn bộ lịch sử trước đó.
\end{quote}

\subsection{Value function và Bellman equation --- bắt buộc}

Cần hiểu ba đại lượng:

\begin{equation}
V^\pi(s),
\end{equation}

là giá trị của trạng thái khi đi theo policy $\pi$.

\begin{equation}
Q^\pi(s,a),
\end{equation}

là giá trị của việc thực hiện hành động $a$ tại trạng thái $s$, sau đó đi theo policy $\pi$.

\begin{equation}
A^\pi(s,a)=Q^\pi(s,a)-V^\pi(s),
\end{equation}

là mức độ hành động $a$ tốt hơn hoặc kém hơn mức trung bình tại trạng thái $s$.

Bellman equation diễn tả ý tưởng:

\begin{equation}
\text{Giá trị hiện tại}
=
\text{phần thưởng trước mắt}
+
\text{giá trị tương lai}.
\end{equation}

Đây là cầu nối trực tiếp dẫn đến Q-Learning.

\subsection{Dynamic Programming --- rất nên học trước Q-Learning}

Nên biết:

\begin{itemize}[leftmargin=1.5em]
    \item Policy evaluation.
    \item Policy improvement.
    \item Policy iteration.
    \item Value iteration.
\end{itemize}

Dynamic Programming giả sử đã biết mô hình môi trường. Q-Learning thì không cần biết trước mô hình, nhưng cách cập nhật Q-Learning bắt nguồn từ Bellman optimality equation.

\subsection{Exploration và exploitation --- bắt buộc}

Agent phải cân bằng giữa:

\begin{itemize}[leftmargin=1.5em]
    \item \textbf{Exploitation}: chọn hành động đang được cho là tốt nhất.
    \item \textbf{Exploration}: thử hành động khác để tìm lựa chọn tốt hơn.
\end{itemize}

Cần hiểu ít nhất:

\begin{itemize}[leftmargin=1.5em]
    \item Random policy.
    \item Greedy policy.
    \item $\epsilon$-greedy.
    \item Cách giảm dần epsilon.
    \item Vì sao chỉ chọn hành động tốt nhất hiện tại có thể khiến agent mắc kẹt.
\end{itemize}

\subsection{Monte Carlo, Temporal Difference và SARSA --- nên biết}

Không nên chỉ học duy nhất Q-Learning. Trình tự hợp lý là:

\begin{enumerate}[leftmargin=1.8em]
    \item Multi-armed bandit.
    \item Monte Carlo prediction.
    \item Temporal Difference, đặc biệt là TD(0).
    \item SARSA.
    \item Q-Learning.
\end{enumerate}

Công thức Q-Learning:

\begin{equation}
Q(s_t,a_t)
\leftarrow
Q(s_t,a_t)
+
\alpha
\left[
r_{t+1}
+
\gamma \max_a Q(s_{t+1},a)
-
Q(s_t,a_t)
\right].
\end{equation}

Trong đó:

\begin{equation}
\underbrace{
r_{t+1}
+
\gamma \max_a Q(s_{t+1},a)
-
Q(s_t,a_t)
}_{\text{TD error}}
\end{equation}

là sai lệch giữa giá trị mục tiêu mới và giá trị Q hiện tại.

\subsection{Python, NumPy và môi trường RL --- cần cho thực hành}

Bạn nên có khả năng:

\begin{itemize}[leftmargin=1.5em]
    \item Viết vòng lặp \texttt{for}, \texttt{while}.
    \item Dùng hàm và class cơ bản.
    \item Làm việc với NumPy array.
    \item Lập chỉ mục ma trận.
    \item Sampling ngẫu nhiên.
    \item Vẽ reward theo episode.
    \item Đọc API của môi trường.
\end{itemize}

Công cụ thực hành phù hợp:

\begin{itemize}[leftmargin=1.5em]
    \item Python.
    \item NumPy.
    \item Matplotlib.
    \item Gymnasium.
    \item PyTorch khi chuyển sang Deep RL.
\end{itemize}

Không nên chỉ học Q-Learning bằng Matlab nếu mục tiêu dài hạn là nghiên cứu hoặc làm dự án AI hiện đại.

\section{Điều kiện để chuyển từ RL cơ bản sang Deep RL}

Sau khi làm được Q-Learning bằng Q-table, bạn mới nên chuyển sang DQN.

Trước DQN cần hiểu:

\begin{itemize}[leftmargin=1.5em]
    \item Function approximation.
    \item Regression với neural network.
    \item Mean squared error.
    \item Gradient descent.
    \item Input state và output Q-values.
    \item Replay buffer.
    \item Mini-batch sampling.
    \item Target network.
    \item Temporal-difference target.
    \item Vì sao huấn luyện neural network trực tiếp trên dữ liệu RL dễ mất ổn định.
\end{itemize}

Bạn đã có phần lớn nền tảng neural network cho đoạn này. Phần còn thiếu chủ yếu là \textbf{kiến thức RL}, không phải kiến thức CNN.

Sau DQN mới nên học:

\begin{itemize}[leftmargin=1.5em]
    \item Policy Gradient.
    \item REINFORCE.
    \item Actor--Critic.
    \item Advantage Actor--Critic.
    \item PPO.
\end{itemize}

\section{Lộ trình phù hợp nhất với tình trạng hiện tại}

\subsection{Giai đoạn 1 --- Bổ sung nền tảng tối thiểu}

Học:

\begin{enumerate}[leftmargin=1.8em]
    \item Supervised, unsupervised và reinforcement learning khác nhau thế nào.
    \item Linear regression, loss MSE và function approximation.
    \item Xác suất, kỳ vọng và sampling.
    \item Khái niệm Markov property.
\end{enumerate}

Không cần chờ hoàn thành KNN, K-Means và SVM.

\subsection{Giai đoạn 2 --- RL dạng bảng}

Học theo thứ tự:

\begin{enumerate}[leftmargin=1.8em]
    \item Agent--environment interaction.
    \item State, action, reward, policy.
    \item Return và discount factor.
    \item MDP.
    \item State-value và action-value.
    \item Bellman equation.
    \item Policy iteration và value iteration.
    \item Monte Carlo.
    \item TD learning.
    \item SARSA.
    \item Q-Learning.
\end{enumerate}

Thực hành:

\begin{itemize}[leftmargin=1.5em]
    \item Multi-armed bandit.
    \item GridWorld.
    \item FrozenLake.
    \item Taxi.
\end{itemize}

\subsection{Giai đoạn 3 --- Deep RL}

Sau khi tự cài đặt được Q-Learning:

\begin{enumerate}[leftmargin=1.8em]
    \item Function approximation.
    \item DQN.
    \item Experience replay.
    \item Target network.
    \item Double DQN.
    \item Policy Gradient.
    \item Actor--Critic.
    \item PPO.
\end{enumerate}

Thực hành ban đầu với CartPole và LunarLander; chưa cần dùng ảnh nên chưa cần CNN.

\section{Checklist sẵn sàng}

Bạn có thể bắt đầu Q-Learning khi trả lời được các câu sau:

\begin{itemize}[leftmargin=1.5em]
    \item State, action và reward khác nhau thế nào?
    \item Policy là gì?
    \item Episode kết thúc khi nào?
    \item Return khác reward tức thời thế nào?
    \item Discount factor $\gamma$ có tác dụng gì?
    \item Kỳ vọng toán học là gì?
    \item $V(s)$ khác $Q(s,a)$ thế nào?
    \item Bellman equation thể hiện ý tưởng gì?
    \item Exploration khác exploitation thế nào?
    \item Learning rate $\alpha$ điều khiển điều gì?
\end{itemize}

\textbf{Đánh giá hiện tại:} bạn đã sẵn sàng bắt đầu phần nhập môn RL, nhưng nên dành một chặng ngắn bổ sung \textbf{xác suất + MDP + Bellman equation} trước hoặc học song song với nội dung Q-Learning. CNN và AlexNet hiện tại đã nhiều hơn mức cần thiết cho RL cơ bản.

\end{document}
